\Article{363}
Summa rei in hoc vertitur, ut haec residua determinentur. Primo quaeratur residuum, quod periodus maxima per productum ex omnibus functionibus primis idoneis dabit. Si hoc productum sit
\[
\equiv x^\lambda-Ax^{\lambda-1}+\text{etc.}
\]
erit residuum hoc $\equiv A$. Huius autem producti forma facile invenitur et ex Cap. \textsc{vi} sequitur esse $A=0$, si $\nu$ per quadratum dividi possit, contra esse $A$ aut $=+1$ aut $=-1$, prout multitudo factorum primorum numeri $\nu$ sit par aut impar.

Iam resolvatur haec periodus maxima in periodos inferiores repraesententurque periodi cuiusvis termini per $x^{kp^\pi u}$, ita ut $k$ in quavis periodo sit numerus
\SourceLeaf{242}
determinatus, pro diversis vero variabilis, $\pi$ et $u$ autem in quavis periodo variabiles, eos autem valores, quos in aliqua periodo habent, etiam in reliquis adipisci possint. Supponatur aliquantisper aliqua functio prima $P$ pro basi sitque residuum, quod periodi $\Sigma x^{p^\pi u}$, $\Sigma x^{k'p^\pi u}$ etc. per eam divisae praebent respective $A,A'$ etc., erit $\Sigma x^{p^\pi u}-A$ per productum ex omnibus functionibus $(P,\rho^u)$ divisibilis. $\Sigma x^{k'p^\pi u}-A'$ per productum ex omnibus functionibus $(P,\rho^{k'u})$ etc. etc. At facile liquet, quantitates $A,A'$ etc. esse radices congruentiae datae. Scilicet sint periodi radicum aequationis $x^\nu=1$ periodis praecedentibus correspondentes radices aequationis $Q=0$, erunt $A,A'$ etc. radices congruentiae $Q\equiv0$. Namque erit
\[
\begin{aligned}
A+A'+\text{etc.}&\equiv\text{summae periodorum},\\
AA+A'A'+\text{etc.}&\equiv\text{summae quadratorum periodorum}
\end{aligned}
\]
etc. etc. Calculus enim prorsus similis erit ei, quem Cap. \textsc{vi} exposuimus, si pro $\rho$ substituatur $x$, quoniam etiam hic poni potest pro $x^\nu$ unitas, uti illic pro $\rho^\nu$.

Inventis radicibus $A,A'$ etc. aliqua pro residuo periodi $\Sigma x^{p^\pi u}$ eligatur et inde reliquarum residua simili modo uti Cap. \textsc{vi} ordinentur. Namque illud etiam hic arbitrio relinquitur, quum functio $P$ sit prorsus hactenus indeterminata. Calculus sequens omnino analogus est ei, quem Cap. \textsc{vi} pertractavimus, singula exponere nimis prolixum nobis foret. Tandem postquam ad $\Sigma x^{p^\pi}$ perventum est, rei summa perfecta est. Namque posito
\[
P\equiv x^m+ax^{m-1}+bx^{m-2}+\text{etc.}
\]
erit $-a\equiv\Sigma x^{p^\pi}$, eodem modo coëfficiens secundus reliquarum functionum $(P,\rho^k)$, habebitur, unde reliqui ipsius $P$ determinari possunt. Saepius evenire potest, ut ad congruentias identicas perveniatur, ex quibus nihil derivari posse videtur. Quomodo huic difficultati obveniri possit, infra monstrabitur.

\Article{364}
Omnia haec per exemplum multo clariora fient. Resolvenda proponitur functio $x^{15}-1$ secundum modulum $17$ in factores. Erit $m=4$ et quoniam productum ex omnibus functionibus elementaribus
\[
\equiv\frac{x^{15}-1\pdot x-1}{x^3-1\pdot x^5-1}=x^8-x^7+x^5-x^4+x^3-x+1
\]

\SourceLeaf{243}
Quare duo tantummodo erunt factores primi quatuor dimensionum $P$ et $P'$. Iam $x,xx,x^4,x^7,x^8,x^{11},x^{13},x^{14}$ in has duas periodos distribuantur
\[
\begin{aligned}
\Sigma x^{17^\alpha}&\equiv x+xx+x^4+x^8,\\
\Sigma x^{7\pdot17^\alpha}&\equiv x^7+x^{11}+x^{13}+x^{14}
\end{aligned}
\]
Sit secundum alteram functionem $P,P'$
\[
\Sigma x^{17^\alpha}\equiv A,\quad\Sigma x^{7\pdot17^\alpha}\equiv A'
\]
eritque
\[
\begin{aligned}
A+A'&\equiv1\\
AA&\equiv\Sigma x^{2\pdot17^\alpha}+\Sigma x^{3\pdot17^\alpha}+\Sigma x^{5\pdot17^\alpha}+\Sigma x^{9\pdot17^\alpha}\\
A'A'&\equiv\Sigma x^{14\pdot17^\alpha}+\Sigma x^{6\pdot17^\alpha}+\Sigma x^{5\pdot17^\alpha}+\Sigma x^{3\pdot17^\alpha}
\end{aligned}
\]
quare
\[
\begin{aligned}
AA+A'A'&\equiv\Sigma x^{17^\alpha}+\Sigma x^{7\pdot17^\alpha}+4\Sigma x^{3\pdot17^\alpha}+2\Sigma x^{5\pdot17^\alpha}\\
&\equiv1-4-4\equiv-7
\end{aligned}
\]
Hinc $A$ et $A'$ erunt radices congruentiae
\[
xx-x+4\equiv0\;(\mathrm{mod.}\,17)
\]
quae sunt $6,12$. Hinc $P$ dividet
\[
x^8+x^4+xx+x-6
\]
eritque
\[
\equiv x^4-6x^3-2xx-12x+1
\]
$P'$ autem erit $\equiv(P,\rho^7)$ eritque
\[
\equiv x^4-12x^3-2xx-6x+1
\]
